\hypersetup{pdftitle={Ten Problems That Teach Calculus: Teacher Pilot Guide}}
\hero{TEACHER PILOT GUIDE \quad / \quad CORE 10}{Experience before equation}{Two classroom pilots for Ten Problems That Teach Calculus.}

\panel{\textbf{Pilot question.} Can students recognize the mathematical job before you name the method? In 06 they decide what to count. In 09 they decide what counts as a good prediction.}

\sub{At a glance}
{\setlength{\tabcolsep}{8pt}\renewcommand{\arraystretch}{1.16}
\begin{tabularx}{\linewidth}{@{}p{.55in}*{2}{>{\raggedright\arraybackslash}X}@{}}\toprule
 & \textbf{06 \quad Cart on a track} & \textbf{09 \quad Kitchen Inquiry Lab}\\\midrule
Time & 35--45 min & 45--55 min live. Practice: 30--40 min.\\
Group & Pairs, then individual transfer & Teams of 2--4, then individual transfer\\
Before & Print 06A and 06B separately. Set out strips and tokens. & Print 09A and 09B separately. Set out cups and probes or use page 4.\\
After & Worked Problem 06 compares the sums with the exact result. & Worked Problem 09 shows the analytic fit and prediction.\\\bottomrule
\end{tabularx}}

\sub{A release rule that protects the discovery}
Give the A page first and listen to students explain the need in their own words. Then release the B page. For 09, log the original predictions before showing later readings. Introduce formal names after students can describe the jobs without them.

\sub{A common observation sheet}
Collect an initial explanation, a team report, and an individual transfer response. Score each criterion: \textbf{0} absent or incorrect, \textbf{1} partial or prompted, \textbf{2} independent and justified.

{\setlength{\tabcolsep}{7pt}\renewcommand{\arraystretch}{1.16}
\begin{tabularx}{\linewidth}{@{}>{\raggedright\arraybackslash}p{1.30in}*{2}{>{\raggedright\arraybackslash}X}@{}}\toprule
Criterion & \textbf{06 evidence} & \textbf{09 evidence}\\\midrule
Recognizes the need & Distinguishes where it ends from how far it rolled. & Distinguishes planning with early readings from checking with later ones.\\
Explains a method & Uses interval widths and explains which contributions keep their signs. & States a rule and assumption, makes a prediction, and compares with new evidence.\\
Transfers independently & Uses the battery data to justify two distinct answers. & Gives a justified humidity prediction and requests a reading that could test it.\\\bottomrule
\end{tabularx}}

\sub{Make one revision after the pilot}
Note where learners first needed a prompt and which misconception survived. Revise that prompt before expanding to the other eight lessons. A correct number alone is not evidence of transfer.
\foot{Teacher guide \enspace / \enspace Setup and assessment}

\clearpage
\hero{TEACHER PILOT GUIDE \quad 06 \quad / \quad DISTANCE AND DISPLACEMENT}{One trip, two reports}{Keep the sign when it answers the question. Drop it when it does not.}

\sub{Run the room}
\entry{0--8 min}{Move and predict.}{Ask pairs to trace a token backward and forward. If both clients get the same reply, ask what a trip that returns to zero would look like to each client.}
\entry{8--23 min}{Build a rule.}{Give 06B. Ask students to compute one row before filling the table. If they add the velocities, ask what unit the client requested and what the interval length changes.}
\entry{23--32 min}{Compare and formalize.}{Compare reports, ask why the sign matters, then name displacement, total distance, and the definite integral. Use the graph to highlight the turn and the two regions.}
\entry{32--42 min}{Check transfer.}{Have students answer the battery question alone, without recommending a procedure. Collect the rule as well as the numbers.}

\sub{What to listen for}
\textbf{Misconception.} The total is just the absolute value at the end. \textbf{Prompt.} If the cart rolls one meter each way, what does taking the absolute value at the end make invisible? \textbf{Strong response.} Track: combine signed changes. Wheels: make each contribution positive before adding, because both legs put meters on the wheels.

\sub{Key and exact check}
06A: A return to the start gives \(0\uom{m}\) on the track report but positive wheel travel. The quarter-second reading means \(-0.20\uom{m}\) for the track and \(0.20\uom{m}\) for the wheels. On 06B, the four signed contributions are \(-0.46875,\,-0.21875,\,+0.28125,\,+1.03125\) m. Their sum is \(0.625\uom{m}\) forward. Their magnitudes sum to \(2.000\uom{m}\).

For the continuous velocity model, the integral of \(v\) is \(2/3\uom{m}\). The integral of \(\lvert v\rvert\) is \(2\uom{m}\). Here the distance comes out exact even with only four rectangles because the two leg errors cancel after taking magnitudes. It is not a general rule. Shorter slices and an explicit split at a direction reversal improve the method. If velocity changes sign within a slice, a single representative value can hide motion.

\panel{\textbf{Individual transfer key.} The three signed battery contributions are \(+6,-3,+2\) C. Net gain is \(5\) C. Total charge moved is \(11\) C. A strong answer explains why the negative interval counts toward meter throughput and uses all three interval lengths.}

\sub{Optional Python follow-through}
Run companion script \texttt{06\_distance\_vs\_displacement.py}. Ask learners to point to the operation that makes one report different and predict what changes when the plot uses narrower slices.
\foot{Teacher guide \enspace / \enspace Cart investigation}

\clearpage
\hero{TEACHER PILOT GUIDE \quad 09 \quad / \quad KITCHEN INQUIRY LAB}{Make a prediction you can test}{Distinguish the readings that shape the rule from those that can challenge it.}

\sub{Set up the live run}
Use a stable cup on a tray, water, a thermometer, and a timer. Record cup material, fill volume, probe placement, room temperature, and whether the cup is covered. Have the teacher handle hot water and supervise the probe, keeping it off the walls and bottom. Test your setup in advance so the drink stays above \(40\tc\) through the early readings and crosses during class. If it doesn't cross, record an incomplete test and use the practice record for a complete one. Use the practice record if live hot-water work is unsuitable.

\sub{Sequence and release}
\entry{Before time 0}{Predict.}{Ask whether temperature will fall the same number of degrees every minute, and what happens near room temperature. Do not give the equation.}
\entry{0, 4, 8 minutes}{Collect early readings.}{Let teams propose a rule. Encourage them to work in degrees above room temperature and compare what gets lost with what remains. Use the first and last readings to fit and the middle one as an early check.}
\entry{After the fit}{Commit and hide.}{Log each team's original prediction and its assumption. Keep later readings out of view until every team commits. For live runs, log every minute near the predicted crossing if feasible, in addition to the planned times.}
\entry{After the crossing}{Test and explain.}{Release 09B and the later data. Compare the original prediction with the observed bracket and interpolated crossing. A refit to the late data is a revision, not fresh validation.}
\entry{Final 6--8 min}{Transfer.}{Use the humidity question individually. Require a model choice, not just an answer.}

\sub{Prompts that do not give the method away}
\textbf{If a team uses degrees per minute:} What would that rule predict if you ran it long enough? \textbf{If it uses every data point to fit:} Which number is still free to disagree with the prediction? \textbf{If there is no crossing:} What range do the observations actually support?

\panel{\textbf{Strong student explanation.} The drink's excess above room temperature shrinks by the same fraction in equal intervals under our rule. We used two early readings to set the rule and later ones to evaluate it. The result is useful only if it meets the client's accuracy requirement.}

\sub{When to name the calculus}
After teams describe why a larger excess means faster cooling, introduce Newton's model, \(dT/dt=-k(T-T_{\mathrm{room}})\), for constant room temperature and \(k\). For students not ready for logarithms, graph repeated fractional steps and interpolate. Offer equations after predictions, not on 09A.
\foot{Teacher guide \enspace / \enspace Cooling investigation}

\clearpage
\hero{TEACHER PILOT GUIDE \quad 09 \quad / \quad DATA AND KEY}{A practice run with a real test}{The readings are synthetic. Release the late values only after predictions are logged.}

\sub{Release cards}
{\setlength{\tabcolsep}{7pt}\renewcommand{\arraystretch}{1.24}
\begin{tabularx}{\linewidth}{@{}l*{7}{>{\centering\arraybackslash}X}@{}}\toprule
Time (min) & 0 & 4 & 8 & 12 & 16 & 20 & 24\\\midrule
Drink \(({}^\circ\mathrm C)\) & 70.0 & 61.0 & 54.0 & 48.7 & 44.0 & 40.3 & 37.0\\
Release & Early & Early & Early & Later & Later & Later & Later\\\bottomrule
\end{tabularx}}
Room temperature is \(22\tc\). Copy only the four late columns onto the board after every team has locked its prediction.

\sub{Two candidate rules and the check}
\textbf{Straight-line baseline.} A drop of 16 degrees in 8 minutes is 2 degrees per minute. That would reach \(40\tc\) at 15 minutes. This misses the later crossing.

\textbf{Constant fractional decline in excess.} From \(70-22=48\) to \(54-22=32\), the 8-minute multiplier is \(32/48=2/3\). This gives
\[
 T(t)=22+48\left(\frac23\right)^{t/8},
 \qquad t_{40}=\frac{8\ln(18/48)}{\ln(2/3)}\approx19.35\uom{min}.
\]
The predicted 4-minute reading is \(61.19\tc\) versus \(61.0\tc\) in the early record. Predicted temperatures at 12, 16, 20, and 24 minutes are \(48.1,\ 43.3,\ 39.4,\ 36.2\tc\). The later readings are all somewhat warmer.

\textbf{Use the later evidence, with its limits.} The crossing is bracketed by 20 and 24 minutes. Assuming linear change between those readings gives \(20+4(40.3-40)/(40.3-37.0)\approx20.36\uom{min}\). The estimated error is \(1.01\) minutes, inside the 2-minute tolerance. The bracket alone does \emph{not} establish that verdict. A crossing after approximately 21.35 minutes would fail. In a fresh trial, add measurements around the crossing. At room temperature \(40\tc\), the ideal model approaches the target without reaching it in finite time from above. In a warmer room it cannot reach it from above.

\panel{\textbf{Individual transfer key.} Assume constant fractional decay of humidity above the expected \(40\%\) equilibrium. The excess falls from 20 to 15 percentage points in 5 minutes, a factor of \(3/4\). Solving \(20(3/4)^{t/5}=5\) gives \(t\approx24.09\) minutes. Request a fresh reading near the prediction or readings just before and after. Other justified models with explicit assumptions and a meaningful test can earn credit.}

\sub{Optional Python follow-through}
Companion script \texttt{09\_coffee\_cooling.py} uses the original textbook values. Agreement between its formula and Euler simulation checks the implementation. It doesn't test a physical cup.
\foot{Teacher guide \enspace / \enspace Key and evidence}
